% Ad Familiares 13.54 --- headnote
% ad-familiares-13-54, between 11 February and April 50 BC, Laodicea

\textit{Cicero to Q. Minucius Thermus, propraetor of Asia, written
from Laodicea between the third day before the Ides of February and
the month of April (between 11 February and April) 50 BC (the
manuscript dateline: \textit{Scr. Laudiceae inter a. d. iii Id.
Febr. et m. Apr. a. 704 (50)}). One of the recommendation-letters
that fill book 13 of \textit{Ad Familiares}: the proconsul of
Cilicia writing to a neighbouring provincial governor on behalf of
the son of his interpreter Marcilius.}

\textit{Cicero has already commended young Marcilius to Thermus once;
this letter circles back to thank him for his generous reception of
the man at Laodicea, and to ask him to go a little further --- to
take care, ``so far as your honour permits,'' that the boy's
mother-in-law not be made the defendant in some pending case. The
elder Marcilius, who has served on Cicero's staff in Cilicia, gets
the closing accolade: \textit{singularis et prope incredibilis fides,
abstinentia modestiaque}.}
